% Part: model-theory
% Chapter: lindstrom
% Section: introduction

\documentclass[../../../include/open-logic-section]{subfiles}

\begin{document}

\olfileid[fr]{mod}{lin}{int}
\olsection{Introduction}

Nous voulons démontrer la caractérisation,
due à Lindström, de la logique du premier
ordre : sous certaines autres conditions,
c'est la logique la plus expressive qui
possède à la fois la propriété de compacité
et la propriété de Löwenheim--Skolem
descendante
(\olref[fol][com][com]{thm:compactness} et
\olref[fol][com][dls]{thm:downward-ls}).
Il nous faut d'abord définir plus
généralement la classe des logiques
auxquelles s'applique ce théorème.
Nous nous limiterons aux langages
\emph{relationnels}, c'est-à-dire
aux langages qui ne contiennent que
des !!{predicate}s et des constantes
individuelles, sans !!{function}s.

\end{document}
